\documentclass[UTF8,a4paper,11pt]{ctexart}
\usepackage[margin=2.4cm]{geometry}
\usepackage{booktabs}
\usepackage{longtable}
\usepackage{array}
\usepackage{amsmath}
\usepackage{hyperref}

\hypersetup{colorlinks=true,linkcolor=blue,urlcolor=blue}
\setlength{\emergencystretch}{3em}
\newcolumntype{L}[1]{>{\raggedright\arraybackslash}p{#1}}

\title{实验数据集与参数设置}
\author{Bridge DL Poly Experiments}
\date{\today}

\begin{document}
\maketitle

\section{实验设置总览}

本项目的实验分为四条主线：低维 NN--多项式几何等价性实验、pilot 距离驱动的 D-optimal 迁移实验、高维迭代 D-optimal 实验，以及最终的 measure-weighted 批量采样实验。所有实验的共同目标是检验 NN 与多项式回归模型之间的几何距离是否能够解释预测性能差距，并进一步作为 D-optimal 采样迁移的诊断量或策略变量。

\begin{longtable}{L{0.22\textwidth}L{0.30\textwidth}L{0.38\textwidth}}
\caption{实验主线及用途}
\label{tab:experiment-lines} \\
\toprule
\textbf{Experiment} & \textbf{Main Scripts} & \textbf{Purpose} \\
\midrule
\endfirsthead
\toprule
\textbf{Experiment} & \textbf{Main Scripts} & \textbf{Purpose} \\
\midrule
\endhead
\bottomrule
\endlastfoot

Low-dimensional geometry experiment
& \texttt{measure\_morala.py}; \texttt{analyze\_results.py}
& Tests whether geometric closeness between NN and polynomial regression predicts predictive-performance closeness. \\

Pilot D-optimal transfer
& \texttt{d\_optimal\_transfer}; \texttt{pilot\_dopt\_transfer}
& Tests whether a small pilot NN--FullPR distance can diagnose when FullPR-based D-optimal sampling transfers to NN distillation. \\

High-dimensional iterative D-optimal
& \texttt{iterative\_highdim\_dopt}; \texttt{paper\_highdim\_iter}
& Tests whether repeatedly upgrading the design set by D-optimal leverage improves a FullPR surrogate of an NN oracle. \\

Measure-weighted batch sampling
& \texttt{measure\_weighted\_sampling}; \texttt{combine\_measure\_weighted\_grand}
& Uses NN--PR/TYPR distances as a transfer gate to mix D-optimal and random sampling. \\

\end{longtable}

\section{Datasets}

\subsection{Low-dimensional datasets}

The low-dimensional geometry experiments use $n=200$ samples and $p=3$ input variables. Inputs are generated before scaling, then both $X$ and $y$ are min--max scaled to $[-1,1]$ after a fixed 75/25 train/test split. The data seed is fixed at 0 in the final grid, while NN initialization seeds vary.

\begin{longtable}{L{0.22\textwidth}L{0.34\textwidth}L{0.34\textwidth}}
\caption{Low-dimensional datasets used in the geometry and pilot-transfer experiments}
\label{tab:lowdim-datasets} \\
\toprule
\textbf{Dataset} & \textbf{Generation / Structure} & \textbf{Role in the Experiments} \\
\midrule
\endfirsthead
\toprule
\textbf{Dataset} & \textbf{Generation / Structure} & \textbf{Role in the Experiments} \\
\midrule
\endhead
\bottomrule
\endlastfoot

\texttt{paper\_poly3}
& Complete degree-3 polynomial with interaction terms and Gaussian noise.
& Well-specified control case when the NN depth matches the polynomial order. \\

\texttt{paper\_poly4}
& Complete degree-4 polynomial with interaction terms and Gaussian noise.
& Strictly under-specified case for the tested FullPR depth range. \\

\texttt{smooth\_nonlinear}
& Fixed smooth non-polynomial target:
$y=\sin(x_1)+x_2^2+x_2x_3+x_3+\epsilon$.
& Smooth non-polynomial benchmark with stable structure. \\

\texttt{smooth\_nonlinear\_rand}
& Random-frequency non-polynomial target:
$y=a\sin(\omega x_1+\phi)+x_2^2+x_2x_3+x_3+\epsilon$.
& Harder non-polynomial benchmark where the sinusoidal component changes by seed. \\

\end{longtable}

\subsection{High-dimensional datasets}

The high-dimensional experiments use $n=10000$ original samples, a 75/25 train/test split, and $p\in\{20,50\}$. Inputs are generated from a standard normal distribution, split into training and test sets, and scaled to $[-1,1]$. The response is also scaled to $[-1,1]$. The noise level is $\sigma=0.03$.

\begin{longtable}{L{0.24\textwidth}L{0.36\textwidth}L{0.30\textwidth}}
\caption{High-dimensional datasets used in the D-optimal and measure-weighted experiments}
\label{tab:highdim-datasets} \\
\toprule
\textbf{Dataset} & \textbf{Generation / Structure} & \textbf{Role in the Experiments} \\
\midrule
\endfirsthead
\toprule
\textbf{Dataset} & \textbf{Generation / Structure} & \textbf{Role in the Experiments} \\
\midrule
\endhead
\bottomrule
\endlastfoot

\texttt{highdim\_poly2}
& Random quadratic polynomial generated from degree-2 polynomial features.
& High-dimensional polynomial reference case. \\

\texttt{highdim\_smooth}
& Quadratic backbone plus moderate smooth nonlinear terms, including sine, cosine, and tanh components.
& Mild nonlinear high-dimensional case. \\

\texttt{highdim\_strong}
& Quadratic backbone plus stronger nonlinear sine, cosine, and tanh terms.
& Boundary case where the NN oracle is less fully represented by degree-2 FullPR. \\

\texttt{highdim\_local}
& Low-amplitude global terms plus localized Gaussian bump and ridge-like structure.
& Tests whether sampling can detect local interior structure. \\

\texttt{highdim\_highfreq}
& Low-amplitude global terms plus high-frequency sine/cosine components.
& Stress case for negative transfer and over-aggressive D-optimal sampling. \\

\end{longtable}

\section{Parameter Settings}

\subsection{Low-dimensional geometry experiment}

\begin{longtable}{L{0.30\textwidth}L{0.58\textwidth}}
\caption{Parameter settings for the low-dimensional geometry experiment}
\label{tab:lowdim-parameters} \\
\toprule
\textbf{Parameter} & \textbf{Setting} \\
\midrule
\endfirsthead
\toprule
\textbf{Parameter} & \textbf{Setting} \\
\midrule
\endhead
\bottomrule
\endlastfoot

Datasets
& \texttt{paper\_poly3}, \texttt{paper\_poly4}, \texttt{smooth\_nonlinear}, \texttt{smooth\_nonlinear\_rand}. \\

Sample size and dimension
& $n=200$, $p=3$. \\

Train/test split
& 75/25 split with fixed \texttt{random\_state=42}. \\

Scaling
& Min--max scaling of both $X$ and $y$ to $[-1,1]$. \\

Activation functions
& \texttt{softplus}, \texttt{tanh}, \texttt{sigmoid}; \texttt{smooth\_nonlinear} is run only with \texttt{softplus} in the final grid. \\

Hidden-layer structures
& $(4)$, $(8,4)$, $(16,8,4)$. \\

Taylor truncation order
& $q=3$. \\

FullPR feature setting
& Both \texttt{include\_special=True} and \texttt{include\_special=False}; special terms include sinusoidal/exponential features. \\

FullPR maximum order
& Equal to the number of NN hidden layers in each configuration. \\

Optimizer and training epochs
& RPROP optimizer; 1500 epochs in the default experiment setting. \\

Repeated runs
& 50 NN initialization seeds per configuration; fixed data seed 0. The final low-dimensional result contains 3000 runs. \\

Main outputs
& \texttt{results\_raw.csv}, cleaned correlations, LTPR failure rates, NN-baseline comparisons, geometry/MSE figures. \\

\end{longtable}

\subsection{Pilot D-optimal transfer experiment}

\begin{longtable}{L{0.30\textwidth}L{0.58\textwidth}}
\caption{Parameter settings for the pilot D-optimal transfer experiment}
\label{tab:pilot-dopt-parameters} \\
\toprule
\textbf{Parameter} & \textbf{Setting} \\
\midrule
\endfirsthead
\toprule
\textbf{Parameter} & \textbf{Setting} \\
\midrule
\endhead
\bottomrule
\endlastfoot

Datasets
& \texttt{paper\_poly3}, \texttt{paper\_poly4}, \texttt{smooth\_nonlinear}, \texttt{smooth\_nonlinear\_rand}. \\

Activations
& \texttt{tanh} for \texttt{paper\_poly3}, \texttt{paper\_poly4}, and \texttt{smooth\_nonlinear\_rand}; \texttt{softplus} for \texttt{smooth\_nonlinear}. \\

Full data size
& $n_{\mathrm{full}}=5000$. \\

NN architecture
& Hidden layers $(16,8,4)$. \\

Training epochs
& 500 epochs in the pilot-transfer runner. \\

Pilot fractions
& 0.005, 0.01, and 0.1 in the final CPU/GPU summary. \\

Candidate and evaluation points
& 3000 candidate points and 1000 evaluation points. \\

Budgets
& 26 and 52 queried NN points in the final combined result. \\

Random baseline
& Same-budget random designs, repeated 30 times. \\

Devices
& CPU and CUDA. \\

Main outputs
& Pilot distance, D-optimal MSE gain versus random median, CPU/GPU comparison tables and figures. \\

\end{longtable}

\subsection{High-dimensional iterative D-optimal experiment}

\begin{longtable}{L{0.30\textwidth}L{0.58\textwidth}}
\caption{Parameter settings for the high-dimensional iterative D-optimal experiment}
\label{tab:highdim-iter-parameters} \\
\toprule
\textbf{Parameter} & \textbf{Setting} \\
\midrule
\endfirsthead
\toprule
\textbf{Parameter} & \textbf{Setting} \\
\midrule
\endhead
\bottomrule
\endlastfoot

Datasets
& \texttt{highdim\_poly2}, \texttt{highdim\_smooth}, \texttt{highdim\_strong}. \\

Original sample size
& $n=10000$; 7500 training-pool points and 2500 evaluation points. \\

Input dimensions
& $p=20$ and $p=50$. \\

Data seeds and NN initialization seeds
& Data seeds 0, 1, 2; NN initialization seeds 0, 1. \\

Initial uniform fractions
& 0.05 and 0.1, corresponding to 375 and 750 initial design points. \\

Iterations and batch size
& 5 D-optimal upgrade iterations; each iteration adds 500 points. Final selected sizes are 2875 and 3250. \\

NN oracle
& Hidden layers $(128,64,32)$, \texttt{tanh} activation, 1000 epochs. \\

FullPR surrogate
& Degree-2 polynomial features with special terms; ridge parameter $\alpha=10^{-5}$. \\

D-optimal approximation
& Batch leverage approximation with ridge $10^{-2}$ and chunk size 512. \\

Random baseline
& 10 same-budget random upgrades at each iteration point. \\

Devices
& CPU and CUDA. \\

Shape subsampling
& Up to 1000 points are used for shape-distance computation. \\

Main outputs
& \texttt{paper\_highdim\_iter\_dopt\_raw.csv}, aggregate statistics, final summary, MSE/gain curves, and final-gain figures. \\

\end{longtable}

\subsection{Measure-weighted batch sampling experiment}

\begin{longtable}{L{0.30\textwidth}L{0.58\textwidth}}
\caption{Parameter settings for the measure-weighted batch sampling experiment}
\label{tab:measure-weighted-parameters} \\
\toprule
\textbf{Parameter} & \textbf{Setting} \\
\midrule
\endfirsthead
\toprule
\textbf{Parameter} & \textbf{Setting} \\
\midrule
\endhead
\bottomrule
\endlastfoot

Datasets
& \texttt{highdim\_poly2}, \texttt{highdim\_smooth}, \texttt{highdim\_strong}, \texttt{highdim\_local}, \texttt{highdim\_highfreq}. \\

Original sample size
& $n=10000$; 7500 training points and 2500 test points. \\

Input dimensions
& $p=20$ and $p=50$. \\

Data seeds and initial seeds
& Data seeds 0, 1, 2, 3, 4; initial sampling seeds 0, 1. \\

Initial pilot fraction
& 0.05 of the training pool, i.e. 375 initial points. A 25\% split of the pilot subset is used as pilot validation. \\

Batch rounds and batch size
& 6 batch updates; each update adds 500 points. Selected sizes are 375, 875, 1375, 1875, 2375, 2875, and 3375. \\

Pilot NN
& Single-hidden-layer NN with hidden size 64, \texttt{tanh} activation, 350 epochs. \\

Main NN/surrogate fitting
& Same hidden size and activation; 350 epochs for NN training. \\

FullPR and Taylor-PR settings
& FullPR degree 2; Taylor order 3; ridge parameter $\alpha=10^{-4}$. \\

Distance computation
& Shape distance uses at most 256 sampled points. The combined transfer distance is $\mathrm{mean}(d_{\mathrm{FPR}},d_{\mathrm{TYPR}})$. \\

Distance normalization
& Distances are normalized using the 0.10 and 0.90 quantiles. \\

D-optimal weight bounds
& \texttt{dopt\_weight} is bounded between 0.30 and 0.95. \\

Compared strategies
& \texttt{Measure-weighted}, \texttt{Random}, \texttt{D-optimal}, and \texttt{Latin hypercube}. \\

D-optimal approximation
& Batch leverage approximation with ridge $10^{-6}$ and chunk size 2048. \\

Device
& CUDA in the grand benchmark commands. \\

Main outputs
& Combined raw/final/distance CSV files, strategy summaries, case summaries, pairwise comparisons, learning curves, gain plots, and weight heatmaps. \\

\end{longtable}

\section{Notes for Reporting}

Across all experiments, positive gain means that the tested design or sampling strategy obtains a lower MSE than the corresponding random baseline under the same budget. The pilot-distance and measure-weighted experiments should therefore be interpreted as transfer-diagnosis experiments: the distance is not claimed to make D-optimal universally optimal, but to indicate when old PR/D-optimal structure is likely to transfer to a new NN model.

\end{document}
